\subsection{Descent of morphisms}

Let X and Y be topological spaces, and let f be a continuous map from X to Y. Let \((Z,p)\) and \((Z',p')\) be X-spaces equipped with descent data relative to f, denoted respectively \(\tau\) and \(\tau'\) . We say that an X-morphism \(\varphi\colon Z\to Z'\) is compatible with the descent data \(\tau\) and \(\tau'\) if one has

\[
\tau^ {\prime} (\varphi (z), x) = \varphi (\tau (z, x))
\]

for all \((z,x)\in\mathbf{Z}\times_{\mathbf{Y}}\mathbf{X}\) . It is equivalent to say that the images under \(\varphi\) of two equivalent points with respect to the relation \(R_{\tau}\) are equivalent with respect to the relation \(R_{\tau^{\prime}}\) . Such a morphism \(\varphi\) defines, by passing to quotients, a continuous map \(\overline{\varphi}:Z/R_{\tau}\to Z^{\prime}/R_{\tau^{\prime}}\) ; it is a morphism of Y-spaces.

Denote by \(\mathcal{C}_{\tau,\tau^{\prime}}(\mathrm{Z};\mathrm{Z}^{\prime})\) the set of X-morphisms from Z to \(Z^{\prime}\) that are compatible with the descent data \(\tau\) and \(\tau^{\prime}\) .

PROPOSITION 3. --- Let X and Y be topological spaces and let \(f\colon X \to Y\) be a continuous map. Let \((Z,p)\) and \((Z',p')\) be X-spaces equipped with descent data relative to f , denoted respectively \(\tau\) and \(\tau'\) . If the descent datum \(\tau'\) is effective, the map \(\varphi \mapsto \overline{\varphi}\) is a bijection from \(\mathcal{C}_{\tau,\tau'}(Z;Z')\) onto \(\mathcal{C}_Y(Z/R_\tau;Z'/R_{\tau'})\) .

For every X-morphism \(\varphi\) from Z into \(Z'\) compatible with the descent data, the map \(\overline{\varphi}\) is a Y-morphism. Conversely, let \(g\colon Z\to Z/R_{\tau}\) and \(g'\colon Z'\to Z'/R_{\tau'}\) be the canonical maps, and let \(\psi\colon Z/R_{\tau}\to Z'/R_{\tau'}\) be a Y-morphism. The maps \(p\colon Z\to X\) and \(\psi\circ g\colon Z\to Z'/R_{\tau'}\) are Y-morphisms. The hypothesis that the descent datum \(\tau'\) is effective means that the diagram

\[
\begin{array}{c} \mathrm {Z^ {\prime}} \xrightarrow {g ^ {\prime}} \mathrm {Z}^ {\prime} / \mathrm {R}_ {\tau^ {\prime}} \\ p ^ {\prime} \Biggl \downarrow \qquad \qquad \qquad \qquad \Biggl \downarrow q ^ {\prime} \\ \mathrm{X} \xrightarrow {f} \mathrm{Y} \end{array}
\]

is a Cartesian square. There therefore exists a unique continuous map \(\varphi\colon Z\to Z'\) such that \(p'\circ\varphi=p\) and \(g'\circ\varphi=\psi\circ g\) . The first equality means that \(\varphi\) is an X-morphism, the second equality means that \(\varphi\) is compatible with the descent data and that \(\overline{\varphi}=\psi\) , whence the proposition.

